\scanpage{36}
L'élément $\alpha\in P(\mathbf C)$ est relié à $f_\infty$ par
\[
\overline\alpha=f_\infty\alpha
\]
donc $f_\infty$ à priori
\[
\left(f_\infty=1\Longleftrightarrow\alpha\in P(\mathbf R)\quad\text{i.e.}\quad\mathbf R(\alpha)=\mathbf R\right)
\]
\deleted{Comme on a la relation $f_\infty\unsure{j_1}f_\infty=\unsure{j_2}$, cette condition implique $j_1=j_2$ donc $j_1=j_2=\frac12 i$, i.e. \deleted{$\unsure{G}$} [si on admet conj. de Hodge] $G^0\xrightarrow[\sim]{\varepsilon}G_m$}

Or on vérifie aussitôt
\[
\boxed{\varepsilon(f_\infty)=-1}
\]
ce qui nous montre qu'on a \unsure{tjrs}
\[
\boxed{f_\infty\ne1\quad\text{i.e.}\quad\mathbf R(\alpha)=\mathbf C}
\qquad\text{(si $M$ contient le motif de Tate)}
\]
Si $M$ ne contient pas le motif de Tate, on voit que $f_\infty=1\Longrightarrow j_1=j_2$ i.e. $j_1=j_2=\frac12 i$, donc en vertu de la conj. de Hodge, il faut que $G^0=\operatorname{Im}\unsure{j_2|j_1}$, donc ou bien $G^0=\{e\}$ \deleted{(motifs de dimension 0, sur $G=\{e\}$ \unsure{ou} $\mathbf Z/2\mathbf Z$, dans ce dernier cas $f_\infty$} \inserted{ici $G=G^0$ car $\unsure{f_\infty}$} \inserted{donc $G=1$} [\uline{NB} si $G$ est disconnecté, on sait que $f_\infty\notin G^0(\mathbf Q)$], ou bien $G=G_m/\mu_n$, $n\geq1$, $M$ est \deleted{\unsure{produit de}} engendré par $\mathbf Q(-n)$\} ; on trouve en $\unsure{=}$ temps que $n$ est \uline{pair}. Donc on trouve

\diagram{\begin{tikzcd}[ampersand replacement=\&,column sep=large,row sep=large]
(f_\infty=1)\arrow[r,Leftrightarrow]
\& (\mathbf R(\alpha)=\mathbf R)\arrow[r,Leftrightarrow]
\& \begin{array}{c}M\text{ est engendré par }\mathbf Q(-2n),\\\text{un quelconque }n\geq0.\end{array}\arrow[d,Leftrightarrow]\\
\&\&G\simeq G_m\text{ ou }G\simeq\{e\}
\end{tikzcd}}
\margin{moyennant conj. de Hodge}

\margin{\fbox{\parbox{0.88\linewidth}{\uline{NB} $f_\infty$ ne peut être \uline{\unsure{central}} que si $G^0=\operatorname{Im}\deleted{\unsure{j}}i$, ce qui implique que $M$ est \deleted{engendré} contenu dans la cat. engendrée par $\mathbf Q$ \unsure{tordu} et par $\mathbf Q(2)$. Il revient au $\unsure{=}$ dire que $G$ est \uline{\unsure{comm.}}}}}
\mspage{35}
